\documentclass[11pt]{article}

\usepackage[margin=1in]{geometry}
\usepackage[T1]{fontenc}
\usepackage{lmodern}
\usepackage{amsmath,amssymb,amsthm}
\usepackage{bm}
\usepackage{booktabs}
\usepackage{graphicx}
\usepackage{siunitx}
\usepackage{tabularx}
\usepackage{array}
\usepackage[round,numbers,sort&compress]{natbib}
\usepackage[colorlinks=true,linkcolor=blue,citecolor=blue,urlcolor=blue]{hyperref}

\title{A Stochastic Groundwater Remediation Problem solved 25 years later with AI\\
\large Introduction}
\author{Sergio Lucero}
\date{\today}

\begin{document}
\maketitle

\section{Introduction}
\subsection{Twenty-five years ago, at Davis}

In 2001, as a doctoral student in Applied Mathematics at the University of
California, Davis, I was handed this problem by my advisor, Roger Wets: a
contaminant plume is moving through an aquifer; a set of wells can pump it out and
treat it; where do you put the wells, and at what rates do you run them, when the
conductivity field that decides where the plume actually goes is known only as a
random field? The committee that examined the result was as interdisciplinary as
the problem itself --- Roger Wets in stochastic programming, Graham Fogg in
hydrogeology, Janko Gravner in probability --- and the question arrived already in
the language of stochastic programming, because that is the language my advisor
spoke. A decision has to be taken before the uncertainty is resolved, the physical
state responds to it, the regulator asks for a probability rather than a
guarantee, and the cost of being wrong is measured in decades of pumping.
Decomposition by scenarios \citep{rockafellar1991}, the L-shaped lineage of
cutting-plane methods \citep{vanslyke1969,benders1962}, and the sample average
approximation \citep{kleywegt2002} were the available machinery; the risk measure
that this thesis ends up using was being written down in the same years
\citep{rockafellar2000}.

The problem had exactly the structure those methods were built for, and it had
one property they could not tolerate: the evaluation of the objective required
solving partial differential equations on a realization of the conductivity
field. Not one solve --- a scenario set of them, for every candidate decision,
inside every iteration of every subproblem. At 2001 prices that was simply not
affordable: a compass sweep over nine well rates, at the scenario counts a tail
statistic needs, was days of computation per iteration, and the search that was
supposed to consume the answers did not exist yet either.

So what I built, and what the 2001 thesis contains, is the forward model: a
deterministic flow-and-transport simulation in which a candidate pump-and-treat
design is evaluated rather than searched for. That is the piece this document
still stands on --- every later chapter evaluates (P) through exactly such a
model --- but it is not (P). The optimization, the scenario ensemble, the risk
measure and any notion of a certificate were named in that thesis and none of them
were computed. The field as a whole did better in one direction and worse in
another: assembling a response matrix from unit-rate solves
\citep{gorelick1983,wagner1995,mayer2002,ahlfeld2000} makes optimization
affordable by replacing the simulator with a linear surrogate, and the risk
measure is then flattened to a variance or the chance constraint to a Gaussian
assumption, so the design is optimized against the surrogate and validated
against the realizations that built it. Those designs were defensible. The
\emph{evidence} for them was thin, and thin in the way that is hardest to notice.

What this document therefore starts from is a problem stated correctly in 2001, a
forward model built in answer to it, and the program itself left open. It returns
to that program --- written again below, in essentially the same form --- and what
changed is not the formulation but the price of an answer, and the standard of
proof applied to it.

\subsection{The decision, and why it is a stochastic program}

Pumping groundwater out of an aquifer, treating it at the surface and returning
it is still the most widely used remedy for a dissolved contaminant plume, and it
is a design problem before it is an engineering one. The physics is classical:
steady Darcian flow through a heterogeneous conductivity field carries a
dissolved species that spreads by dispersion, and the design objective is to keep
the concentration at a compliance surface below a regulatory limit while pumping
as little water as possible for as short a time as possible. What makes the
problem hard is not the partial differential equation. It is that the
conductivity field which determines where the plume goes is known only through a
statistical model fitted to a handful of boreholes, so the object that has to be
optimized is not a number computed from a field but an expectation over fields,
and the constraint that matters is a statement about a distribution.

Written out, the problem this thesis solves is

\begin{equation}
\begin{aligned}
(\mathrm{P})\qquad
\min_{\mathbf q\in\mathcal Q}\quad
& \mathbb{E}\bigl[J(\mathbf q,\xi)\bigr]
  \;+\; \rho_g\Bigl[\operatorname{CVaR}_{\alpha}\bigl(G(\mathbf q,\xi)\bigr)\Bigr]_+^{2}
  \;+\; \rho_h\Bigl[\operatorname{CVaR}_{\alpha_h}\bigl(-H(\mathbf q,\xi)\bigr)\Bigr]_+^{2}
\\[0.5em]
\text{subject to}\quad
& \nabla\cdot\bigl(K(\mathbf x,\xi)\,\nabla h\bigr)
   = \sum_{i=1}^{n_w} q_i\,\delta(\mathbf x-\mathbf x_i),
  \qquad \mathbf v = -\frac{K(\mathbf x,\xi)}{\phi}\,\nabla h,
\\
& \phi R\,\partial_t c
   + \nabla\cdot\bigl(\mathbf v\,c - \mathbf D\,\nabla c\bigr)
   = s(\mathbf x,t),
\\
& G(\mathbf q,\xi) = \max_{t\ge t_0} c(\mathbf x_c,t;\mathbf q,\xi) - c_{\lim},
  \qquad H(\mathbf q,\xi) = \min_{j}\bigl(h(\mathbf x_j,t;\mathbf q,\xi) - h_{\min}\bigr),
\\
& \mathbb{P}\bigl(G(\mathbf q,\xi)\le 0\bigr)\;\ge\;1-\alpha,
  \qquad \alpha = 0.05,
\\
& \mathcal Q = \Bigl\{\mathbf q:\; 0\le q_i\le \bar q,\;\;
  \textstyle\sum_{i=1}^{n_w} q_i \le \bar Q\Bigr\},
\end{aligned}
\label{eq:int-problem}
\end{equation}

in which the decision $\mathbf q$ is the vector of pumping rates at
$n_w$ candidate wells; the uncertainty $\xi$ is the log-normal conductivity field
$K(\mathbf x,\xi)$; the state $(h,c)$ is the pair of random fields solving the
Darcy and transport equations; $J$ is the operating cost of the remediation
(energy for pumping, treatment of the extracted water, residual mass left in the
aquifer, and the time until the plume is contained); $G\ge 0$ is the amount by
which the compliance-point concentration exceeds its limit $c_{\lim}$ over the
late part of the horizon; and $H\ge 0$ is the head-feasibility margin at the
drawdown observation points, so that $H<0$ means the wells are drawing the water
table down further than the permit allows. The chance constraint is the
regulatory requirement; the two penalty terms implement it --- and a second,
stricter one on drawdown --- as an exact penalization through the
Rockafellar--Uryasev representation of conditional value-at-risk, which is what
makes the problem differentiable in the risk argument and finite-dimensional in
$\eta$. Chapter~\ref{chap:optimization} shows that this substitution is exact for
$\rho$ large enough, and it is the reason the policies of
Chapter~\ref{chap:results_2d} sit \emph{on} the constraint rather than inside it.

\begin{table}[t]
  \centering
  \small
  \begin{tabularx}{\textwidth}{@{}l l X@{}}
    \toprule
    Object & Symbol & What it is here \\
    \midrule
    decision & $\mathbf q\in\mathbb{R}^{n_w}$ &
      pumping rate of each candidate well; $n_w=9$ on a hexagonal lattice,
      $0\le q_i\le 50$\,m$^3$/d, $\sum_i q_i\le 200$\,m$^3$/d \\
    uncertainty & $\xi$ &
      log-normal $K$ with $\sigma=0.65$ and correlation length $15$\,m,
      truncated at $1\,529$ Karhunen--Lo\`eve modes (95\% of the variance) \\
    state & $(h,c)$ &
      heads and concentrations from the two equations in (P), on
      $160\times96$ cells ($15\,360$ unknowns), $100$ implicit time steps \\
    source & $s$ &
      $10$\,m$^3$/d of water at $100$\,mg/L for $3\,650$\,d, at one point \\
    compliance & $c_{\lim}$ &
      $1.0$\,mg/L at a single compliance point down-gradient; $G$ is the
      overshoot of its maximum over the late horizon \\
    permit & $h_{\min}$ &
      $-10$\,m of drawdown at the observation points; $H$ is the margin \\
    risk & $\alpha,\alpha_h$ &
      $5\%$ for compliance, $1\%$ for drawdown, both as CVaR \\
    cost & $J$ &
      pumping energy, treatment, residual mass and clean-up time; weights in
      Chapter~\ref{chap:optimization} \\
    \bottomrule
  \end{tabularx}
  \caption{What is decided, what is uncertain and what is computed, for the
    two-dimensional benchmark that carries most of this thesis. The
    three-dimensional model of Chapter~\ref{chap:results_3d} changes the geometry,
    not the structure of (P).}
  \label{tab:int-objects}
\end{table}

Two features of (P) are worth naming, because the thesis is organized around
them. First, the constraint set is defined by \emph{stochastic} partial
differential equations: the coefficient $K$ is random, so the state is a random
field, and the constraint functions $G$ and $H$ are nonlinear functionals of that
state. A scenario is therefore not a row of data but a PDE solve. Second, the
objective is an expectation of a quantity that is itself the output of a
simulator, and the decision appears in it only through the solution --- there is
no closed form, no adjoint that the present search uses
(Chapter~\ref{chap:results_3d}), and no convexity to exploit. Everything
expensive in what follows is a consequence of these two sentences.

\subsection{What was missing then, and what changed}

Three things had to change before (P) could be solved as written rather than
approximated into something else. Only one of them is hardware.

\paragraph{The forward solve.} A compass sweep at the benchmark geometry evaluates
eighteen candidate rate vectors on forty scenarios at once, which on one modern
accelerator takes seconds when the candidates are batched into a single
linear-algebra call and the Krylov iteration is compiled
(Chapter~\ref{chap:results_2d}). Per-solve, the hardware of the period was two to
three orders of magnitude slower; more importantly, nothing in that period
evaluated a whole candidate set in one call, so the amortization that makes a
derivative-free search affordable did not exist. An iteration of this search then
cost what the entire campaign costs now, which is why the searches of that era
were over response-matrix surrogates and not over the simulator.

\paragraph{The surrogate.} Where a forward solve is still too expensive --- and in
three dimensions it is, by a factor that grows like the square of the cell count
(Chapter~\ref{chap:results_3d}) --- a learned surrogate can stand in, but only if
it has been verified against an exact solution on a problem where the exact
solution exists. That is what the one-dimensional sandbox of
Chapter~\ref{chap:formulation} is for: the advection--dispersion equation has the
Ogata--Banks solution \citep{ogata1961solution} in uniform flow, so an
architecture can be scored against truth instead of against a previous run of
itself. This was not available in 2001 either, and neither was the
evidence that it needs to be done carefully: networks that fit the training
distribution perfectly can fail to be operators at all
\citep{krishnapriyan2021failure,wang2021gradient}.

\paragraph{The standard of proof.} The third change is free. The statistics used
here --- a Wilson score interval \citep{wilson1927probable} on a fresh draw of
scenarios, a standard error on the tail, and the refusal to rank policies by a
statistic the sample cannot resolve --- were all available in 2000 and cost
nothing but the decision to spend a validation draw instead of reusing the
training one. The finding that the tail statistic does not transport while the
reliability does (Chapter~\ref{chap:results_2d}) is a statement about how many
scenarios a tail needs, not about how fast a computer is; it could have been
measured then. That it was not is the honest answer to the question of why this
thesis exists.

\subsection{Approach: verification first}

The thesis is organized as three layers of increasing physical complexity, each
carrying its own verification rather than deferring it to the end.

\begin{enumerate}
  \item \textbf{A one-dimensional sandbox with a closed-form answer}
        (Chapter~\ref{chap:formulation}). The advection--dispersion equation has the
        Ogata--Banks solution \citep{ogata1961solution} for a step source in a
        uniform flow, which makes it the right place to test a surrogate
        architecture: the reference is exact, cheap, and covers the parameter
        range. Physics-informed neural networks \citep{raissi2019pinn} and
        Fourier neural operators \citep{li2020fourier,li2021fno} are trained
        there and scored against the analytical solution, with the seed spread
        reported rather than the best seed.
  \item \textbf{A two-dimensional finite-volume model and the stochastic program}
        (Chapters~\ref{chap:formulation} and~\ref{chap:optimization}). Darcy flow
        with partially penetrating wells and advection--dispersion transport are
        discretized on the grid of Table~\ref{tab:int-objects} with a
        Karhunen--Lo\`eve truncation of the log-normal conductivity. The chance
        constraint enters as the CVaR penalty of (P) in a sample-average
        approximation \citep{rockafellar2000,shapiro2009}, and the resulting
        non-convex program is solved by a derivative-free pattern search in which
        an entire compass sweep of candidate rate vectors is evaluated in one
        batched linear solve.
  \item \textbf{An out-of-sample certificate} (Chapter~\ref{chap:results_2d}). A policy is
        not reported until it has been re-solved on scenarios no run has seen, on
        an independent implementation, with reliability quoted as a Wilson
        interval \citep{wilson1927probable} and the tail quoted with its own
        standard error. This layer is where the thesis's methodological claim
        lives, and it is where the strongest negative result appears.
\end{enumerate}

Dimensional fidelity is checked against the same discipline in a third setting: a
three-dimensional finite-volume solver with partially penetrating wells, whose
policy at the paper resolution ($32\times16\times8$) reproduces to the reported
cost on machines as different as a laptop CPU and a rented GPU, with a mass
ledger closing to machine precision. What the 3D model does \emph{not} have is an
analytical gate of its own, and that gap is stated as a limitation rather than
papered over.

\subsection{Two numbers that are easy to get wrong}

It is worth being concrete about the second question, because it motivates every
verification device in the chapters that follow.

\paragraph{A surrogate can be worse than nothing, and the metric will not say so.}
A Fourier neural operator trained on a family of one-dimensional
advection--dispersion problems reported a relative $L^2$ error of $0.973$. The
number had been recorded as a mediocre result. It is not mediocre: relative error
is normalized by the norm of the target, so a predictor that outputs zero
everywhere scores exactly $1.0$. The operator was worse than predicting nothing.
The diagnosis was structural rather than numerical --- the network consumed two
spatially constant parameter channels, applied a pointwise lift, and then
multiplied Fourier coefficients, which can only ever return a constant field ---
and no amount of training would have fixed it. A single trivial-predictor
baseline would have caught this on the first run.

\paragraph{A tail statistic describes the sample that produced it.} The certified
two-dimensional policy of Chapter~\ref{chap:results_2d} reports
$\mathrm{CVaR}_{5\%}(G) = +0.005$ on the hundred scenarios it was optimized on.
Re-solving the same policy on $400$ scenarios drawn from a seed no run had ever
touched gives $+0.616$. Both numbers are correct; they answer different questions,
and only the second says anything about an aquifer. Section~\ref{chap:results_2d}
takes this seriously enough to measure it: across three disjoint draws the tail
takes the values $+0.616$, $+0.254$ and $+0.334$, and those three differ by less
than their own standard errors, because a $5\%$ tail of $400$ draws is an average
over twenty of the most variable realizations in the set and nothing more. The
reliability over the same draws is $0.978$, $0.975$ and $0.978$. One of these
statistics is a property of the policy; the other is a property of the sample.

\subsection{Contributions}

\begin{enumerate}
  \item \textbf{A verification protocol for learned PDE surrogates that uses seed
        sweeps, not best runs.} The baseline PINN on the one-dimensional sandbox
        has a median full-grid relative error of $1.61\times10^{-2}$ over five
        seeds but a range of $1.17\times10^{-2}$ to $8.14\times10^{-2}$; the
        coordinate-augmented FNO reaches a median of $1.99\times10^{-2}$ at 250
        epochs over five seeds, $11\times$ better than a mean-field predictor,
        against a best-seed $1.46\times10^{-2}$. Reporting the best of five
        would have overstated both.
  \item \textbf{A batched finite-volume Darcy/transport solver with a parity-tested
        GPU backend.} A NumPy reference implementation and a torch backend agree
        to $10^{-12}$--$10^{-10}$ and select the same policy end to end, so the
        speedup is not bought with a different model. Batching a full compass
        sweep is worth $47\times$ over the NumPy reference on a T4 at
        $80\times40$; compiling the Krylov iteration into a single call adds a
        further $2.2$--$2.8\times$ on an A100 at $160\times96$.
  \item \textbf{A chance-constrained formulation of pump-and-treat design solved
        to certification.} Two independent runs, differing only in floating-point
        precision, produce policies $1.2\%$ apart in cost on different well
        subsets ($J = 428\,592$ and $433\,945$), both sitting on the constraint
        as the mathematics requires, in $1.67$\,h and $4.14$\,h on one rented GPU.
  \item \textbf{An empirical result about which statistics of a stochastic
        program transport.} On $400$ unseen scenarios a policy selected from only
        sixteen training scenarios reports reliability $1.000$ in-run and fails
        out of sample ($0.912$, Wilson interval $[0.881,0.936]$), while the
        accepted policy holds at $0.978$ $[0.958,0.988]$. Reliability with its
        interval transports across draws and backends; the CVaR tail does not,
        and ranking policies by tail value at affordable sample sizes is not
        supportable.
  \item \textbf{A three-dimensional reference solver, and where the method
        stops working.} The $32\times16\times8$ preset with $1\,419$ modes yields
        test cost $547\,665$, reliability $0.969$, cleanup time $110.6$\,d and a
        mass balance error of $10^{-15}$, reproduced identically on three
        machines. Doubling the resolution to $8\,000$ cells at $24$ training
        scenarios fails its own reliability gate ($0.875$, and $0.930$ $[0.872,
        0.963]$ on fresh draws) for the reason the 2D campaign already
        established --- at $\alpha M_{\text{train}}=1.2$ the tail is one
        scenario --- and the cost accounting that follows ($8$--$9$\,h per
        converged start) is what makes the algorithm, not the solver, the
        binding constraint: the adjoint exists and the search does not use it, and
        the adaptive design in which progressive hedging would earn its keep is
        formulated and checked but not run at benchmark scale.
\end{enumerate}

\subsection{How to read this thesis}

Chapter~\ref{chap:related_work} places the work in the literature on subsurface simulation,
learned surrogates, optimization under uncertainty and verification practice.
Chapter~\ref{chap:formulation} fixes the governing equations, the dimensionless groups
and the two surrogate architectures, and reports their verification against the
analytical solution. Chapter~\ref{chap:optimization} builds the stochastic program:
cost model, risk measure, sample average approximation, and the derivative-free
search with its batched evaluation. Chapter~\ref{chap:results_2d} reports the
two-dimensional results, the certification protocol, and the calibration
experiment that delimits what a tail statistic can support.
Chapter~\ref{chap:results_3d} repeats the protocol in three dimensions, where it
holds at the paper resolution and fails at double the resolution, and follows the
failure to the cost of a forward solve rather than to the solver itself.
Chapter~\ref{chap:conclusions} collects the conclusions, the limitations of the present
evidence, and the work that would be needed to remove them --- of which the most
substantial is the adaptive design in which the decomposition methods of the
Davis years would finally have something to decompose.

\begin{thebibliography}{99}

\bibitem{ogata1961solution}
A. Ogata and R. B. Banks,
\textit{A solution of the differential equation of longitudinal dispersion in
porous media}, U.S. Geological Survey Professional Paper 411-A, 1961.

\bibitem[Raissi et~al.(2019)]{raissi2019pinn}
M.~Raissi, P.~Perdikaris, and G.~E. Karniadakis.
Physics-informed neural networks: A deep learning framework for solving forward
and inverse problems involving nonlinear partial differential equations.
\textit{Journal of Computational Physics}, 378:686--707, 2019.

\bibitem{li2020fourier}
Z. Li, N. Kovachki, K. Azizzadenesheli, B. Liu, K. Bhattacharya, A. Stuart,
A. Anandkumar, \textit{Fourier Neural Operator for Parametric Partial
Differential Equations}, ICLR, 2021.

\bibitem[Li et~al.(2021)]{li2021fno}
Z.~Li, N.~Kovachki, K.~Azizzadenesheli, B.~Liu, K.~Bhattacharya, A.~Stuart,
and A.~Anandkumar.
Fourier neural operator for parametric partial differential equations.
\textit{International Conference on Learning Representations (ICLR)}, 2021.

\bibitem[Rockafellar and Uryasev(2000)]{rockafellar2000}
R.~T. Rockafellar and S.~Uryasev.
Optimization of conditional value-at-risk.
\textit{Journal of Risk}, 2(3):21--41, 2000.

\bibitem[Shapiro et~al.(2009)]{shapiro2009}
A.~Shapiro, D.~Dentcheva, and A.~Ruszczy\'nski.
\textit{Lectures on Stochastic Programming: Modeling and Theory}.
SIAM, Philadelphia, 2009.

\bibitem[Wilson(1927)]{wilson1927probable}
E.~B. Wilson.
\newblock Probable inference, the law of succession, and statistical inference.
\newblock \emph{Journal of the American Statistical Association},
  22(158):209--212, 1927.

\end{thebibliography}

\end{document}
